\subsection{When to use composite training}

The trade-off has a clean structure:

\begin{itemize}
  \item \textbf{Strong yes} when (a) data is constrained
        ($\lesssim 1{,}000$ problems on a math substrate), (b) the
        deployment scenario can exploit mode-switching inference, and
        (c) the diffusion-axis performance matters.
  \item \textbf{Mild yes} when (a) data is abundant but the
        diffusion-axis advantage outweighs the small AR tax for the
        application, and (b) OOD prose generalisation is not a
        priority.
  \item \textbf{No} when (a) the deployment is pure-AR (no
        mode-switching), (b) the only metric is AR-mode held-out
        perplexity, or (c) OOD prose generalisation matters.
\end{itemize}

\subsection{Why the crossover happens}

We do not provide a formal account of \emph{why} the crossover sits
near $\sim 1{,}000$ problems on GSM8K. Two non-exclusive hypotheses:

\begin{enumerate}
  \item \emph{Regularisation budget}. The diffusion-mode mask-fill
        objective acts as an implicit regulariser on the shared
        backbone. In the data-scarce regime, the AR-only model
        overfits to specific token sequences in GSM8K-train, while
        composite generalises better via the denoising objective. As
        data grows, the AR-only baseline gets the regularisation it
        needs from the data itself, and composite's diffusion
        objective becomes a tax rather than an aid.
  \item \emph{Capacity allocation}. The shared backbone has finite
        capacity. With scarce AR signal, it cannot afford to specialise
        in next-token prediction and the bidirectional/diffusion
        component provides complementary signal. With abundant AR
        signal, specialisation in AR pays more than the diffusion
        signal costs.
\end{enumerate}

\subsection{Relationship to published work}

DiFFPO~\citep{wei2025diffpo} reports discrete-diffusion outperforming
AR in data-constrained settings via RL fine-tuning. Our result is a
\emph{from-scratch joint-training} parallel: same regime (data
scarcity), same direction (diffusion-flavoured training wins), no RL.
This is the first evidence we are aware of that the DiFFPO regime
effect manifests in pre-training as well, at small scale and without
RL.

The compute-matched control closes a gap in
BD3-LMs~\citep{nie2025bd3} and Transfusion~\citep{zhou2024transfusion},
neither of which explicitly compares against a compute-matched
pure-diffusion baseline. Our result is small-scale (60M--300M), but
the apples-to-apples comparison is cleaner than in works that focus
on absolute scale.

\subsection{What this study does not claim}

\begin{itemize}
  \item That composite training is universally better than separated
        training. The trade depends on data regime and which axis
        matters.
  \item That the crossover at $\sim 1$k problems generalises beyond
        the GSM8K substrate. A different substrate would produce a
        different curve.
  \item That mode-switching reliably wins downstream tasks. We have
        $n=3$ data with high per-seed variance; the $n=8$ expansion
        tightens this but does not eliminate it.
  \item That toy-scale results scale to $1$B+ parameters. The trade
        direction may shift with model scale; we cannot test.
\end{itemize}
